\begin{afterword}
\summaryhead

The post separates morality (``Killing people is wrong'') from metaethics (what makes it wrong, or what ``wrong'' means). People agree far more on the first than on the second, yet many insist that if their own metaethic is false, morality falls apart. Since their views conflict, many of them will have to change their metaethics, and the author holds that this is nearly impossible while the fear of losing morality is unanswered.

So the post collects the ``lines of retreat'' set up in earlier posts. Whoever obeys God or any other moral authority still chose to obey, and could save a child from the train tracks without it. Explained things are no less valuable. Moral arguments need not move every possible mind. Reflective loops are not circular logic. Evolution's causal role in morality is not a justifying role. Morality computed in the mind need not be about the mind. Metaethics matters, the author says, because misunderstanding where morality comes from harms many people. The post ends on a joke that withholds ``the real reason.''

\responsehead

The post is an index. It gathers arguments from earlier posts and presents each as a line of retreat, a way to keep one's moral commitments while one's metaethics changes. It says openly that it sets these up before the author's own metaethic is presented, so it asks to be read as reassurance rather than as a theory.

As reassurance, its central argument is sound and old. Whoever follows God, or any other moral authority, decided to follow it, and could act rightly without it. Kant made a closely related point: even the Holy One of the Gospels must first be compared with our ideal of moral perfection. The post respects the argument's reach. It offers it instead of attacking belief in God, and the dialogue it links to, ``Is Morality Given?'', states the limit: that the decision to be moral is mine does not show that morality consists of my preferences. The argument settles who is responsible, not what makes an act right.

Some of the positions the post touches can be placed. It does not take the error theorist's side (``not that no morality exists''), and even for error theorists, rejecting moral thought is compatible with keeping its language in a new sense or replacing it, which fits the post's advice that practice can outlast a theory. The worry that morality computed in the mind makes whatever one thinks right is a standard objection to simple subjectivism, which the Stanford Encyclopedia calls the Missing Fallibility Problem. The post's answer, a calculator that computes 2 + 3 rather than its own output, returns in ``Morality as Fixed Computation.''

One claim is given without an example: ``Poor metaethics forms part of the teachings of many a cult, including the big ones.'' No such teaching is named.

In short: an index of reassurances, offered openly as such, whose central argument is old and used within its limits.
\end{afterword}
